\chapter{矩阵}\label{ch:b1:matrices}

矩阵是用坐标写出的线性映射。本章建立一套“词典”：
复合对应矩阵乘积，双射性对应可逆性，
基变换对应相似变换 ---
以及算法部分：初等行变换，秩与逆矩阵的计算。
矩阵在高中分册中已初步接触，
本章将其建立在
\cref{ch:b1:vspaces,ch:b1:findim,ch:b1:linmaps}的理论之上。

\section{矩阵与线性映射}

\begin{definition}\label{def:b1:matrices:def}
$\mathcal{M}_{n,p}(K)$ 是由标量组成的 $n \times p$ 阵列 $A
= (a_{ij})$ 所构成的向量空间（$i$ 为行，$j$ 为列），其维数为 $np$。
（基为 $E_{ij}$，即仅有单个 $1$ 的矩阵）。给定基
$\mathcal{B} = (e_1, \dots, e_p)$（$E$ 的基）与 $\mathcal{C}$（$F$ 的基，$\dim F = n$），
则 \emph{$u \in \mathcal{L}(E, F)$ 的矩阵}为这样的阵列：
其第 $j$ 列由 $u(e_j)$ 在
$\mathcal{C}$ 中的坐标组成：
\[
\operatorname{Mat}_{\mathcal{B},\mathcal{C}}(u) = (a_{ij}),
\qquad u(e_j) = \sum_{i=1}^{n} a_{ij}\, f_i .
\]
映射 $u \mapsto \operatorname{Mat}_{\mathcal{B},\mathcal{C}}(u)$
是从 $\mathcal{L}(E, F)$ 到
$\mathcal{M}_{n,p}(K)$ 上的同构（\cref{prop:b1:linmaps:basis}：线性映射
恰是对各 $e_j$ 像的任意指定）。
\end{definition}

\begin{example}[导数的矩阵表示]
\label{ex:b1:matrices:derivative}
设 $D(P) = P'$，它定义在 $\R_3[X]$ 上。在单项式基 $(1, X, X^2,
X^3)$ 下：$D(1) = 0$，$D(X) = 1$，$D(X^2) = 2X$，$D(X^3) = 3X^2$，于是
\[
\operatorname{Mat}(D) =
\begin{pmatrix}
0 & 1 & 0 & 0\\
0 & 0 & 2 & 0\\
0 & 0 & 0 & 3\\
0 & 0 & 0 & 0
\end{pmatrix}.
\]
在\emph{除幂}基 $\bigl(1,\ X,\ \frac{X^2}2,\
\frac{X^3}6\bigr)$ 中，每个基向量都映到其前一个
基向量（$D\bigl(\frac{X^k}{k!}\bigr) = \frac{X^{k-1}}{(k-1)!}$），且
矩阵化为纯移位形式：上对角线上全为 1，其余位置为 0。两点启示：
矩阵属于\emph{二元组}
（映射，基），而非映射本身；好的基能使结构
一目了然 --- 移位形式立刻说明
$D^4 = 0$ 在 $\R_3[X]$ 上成立，矩阵的每个幂
都把由 1 构成的对角线再向外推一格。
\end{example}

\begin{definition}[乘积]\label{def:b1:matrices:product}
设 $A \in \mathcal{M}_{n,p}$，$B \in \mathcal{M}_{p,q}$，则
\[
(AB)_{ik} = \sum_{j=1}^{p} a_{ij}\, b_{jk}
\qquad (1 \leq i \leq n,\ 1 \leq k \leq q).
\]
这恰好是复合映射的矩阵：
$\operatorname{Mat}(v \circ u) = \operatorname{Mat}(v)\,
\operatorname{Mat}(u)$（中间的基须相互匹配）。类似地，
若 $X$ 为 $x$ 的坐标列，则 $u(x)$ 的坐标列为
$AX$。
\end{definition}

\begin{proof}[复合公式的证明]
\[
v(u(e_k)) = v\Bigl(\sum_j b_{jk} f_j\Bigr) = \sum_j b_{jk}\, v(f_j)
= \sum_j b_{jk} \sum_i a_{ij}\, g_i
= \sum_i \Bigl(\sum_j a_{ij} b_{jk}\Bigr) g_i . \qedhere
\]
\end{proof}

\begin{proposition}[代数 $\mathcal{M}_n(K)$]\label{prop:b1:matrices:ring}
方阵全体 $\mathcal{M}_n(K)$ 构成一个环（当 $n
\geq 2$ 时不可交换），单位元为 $I_n$；其可逆元群为
\emph{一般线性群}\index{一般线性群}
$GL_n(K)$，与双射自同态相对应。对 $A, B \in
\mathcal{M}_n(K)$，有：
\[
AB = I_n \implies A \in GL_n(K) \text{ 且 } B = A^{-1}
\]
（单侧逆即双侧逆，由
\cref{cor:b1:linmaps:samedim}可知）。
\end{proposition}

\begin{proof}
环公理经由 $\mathcal{L}(E)$ 上的
同构（见\cref{def:b1:matrices:def}）：它把复合化为乘积，
把和化为和，因此结合律、分配律以及 $I_n$ 的作用
都继承自映射的相应事实，
无需逐项验证。
不可交换性：$E_{12}E_{21} = E_{11} \neq E_{22}
= E_{21}E_{12}$。若 $AB = I_n$，则 $a$ 为 $A$ 所对应的自同态，
满足 $a \circ b = \mathrm{id}$，故 $a$ 为满射（$x =
a(b(x))$ 给出每个 $x$ 的一个原像），故在
有限维下为双射（见\cref{cor:b1:linmaps:samedim}）。在 $a
\circ b = \mathrm{id}$ 的左侧复合 $a^{-1}$，得 $b =
a^{-1}$，于是 $b\circ a = \mathrm{id}$ 亦成立：单侧逆
一直都是双侧的 --- 这一便利
仅限于有限维情形。
\end{proof}

\begin{definition}[转置；迹]\label{def:b1:matrices:transpose}
$A = (a_{ij}) \in \mathcal{M}_{n,p}$ 的\emph{转置}为 $A^{\mathsf T}
= (a_{ji}) \in \mathcal{M}_{p,n}$，且满足 $(AB)^{\mathsf T} = B^{\mathsf T}
A^{\mathsf T}$，$(A^{\mathsf T})^{\mathsf T} = A$。方阵的
\emph{迹}\index{迹}记为 $\operatorname{tr} A
= \sum_i a_{ii}$，它是线性的，并且
\[
\operatorname{tr}(AB) = \operatorname{tr}(BA)
\qquad (A \in \mathcal{M}_{n,p},\ B \in \mathcal{M}_{p,n}).
\]
\end{definition}

\begin{proof}[迹恒等式的证明]
$\operatorname{tr}(AB) = \sum_i \sum_j a_{ij} b_{ji}$，且
$\operatorname{tr}(BA) = \sum_j \sum_i b_{ji} a_{ij}$：二者为同一个
二重和。
\end{proof}

\begin{example}[迹的应用]
\label{ex:b1:matrices:tracework}
\cref{ch:b1:linmaps}中的投影，映到
$\operatorname{Vect}(1,1)$ 上、沿 $\operatorname{Vect}(0,1)$，
$p(x, y) = (x, x)$，其矩阵为 $A = \begin{pmatrix} 1 & 0\\ 1 &
0\end{pmatrix}$（标准基下），事实上 $A^2 = A$，且
\[
\operatorname{tr} A = 1 = \operatorname{rk} A ,
\]
这印证了\cref{exo:b1:matrices:8}：对于幂等矩阵，迹
\emph{计量}像的维数，无论矩阵写在哪一组
基下。不变性的机理在于恒等式
$\operatorname{tr}(AB) = \operatorname{tr}(BA)$：
\[
\operatorname{tr}\bigl(P^{-1}(AP)\bigr) =
\operatorname{tr}\bigl((AP)P^{-1}\bigr) = \operatorname{tr} A ,
\]
故与 $A$ 相似的所有矩阵共享同一迹 --- 这是
自同态的第一个\emph{数值不变量}，
行列式将在\cref{ch:b1:det}中与之并列（即下文
周末问题中的数对 $(s, p)$）。
\end{example}

\begin{example}[对称加反对称]
\label{ex:b1:matrices:symsplit}
称 $A$ 为\emph{对称矩阵}，当 $A^{\mathsf T} = A$ 时，
称为\emph{反对称矩阵}，当 $A^{\mathsf T} = -A$ 时。任一方阵都
唯一地分解为两者之和：
\[
A = \underbrace{\frac{A + A^{\mathsf T}}{2}}_{\text{对称}}
+ \underbrace{\frac{A - A^{\mathsf T}}{2}}_{\text{反对称}},
\]
既对称又反对称的矩阵为零（$A = -A$）：两个集合
是 $\mathcal{M}_n(K)$ 的补子空间 --- 恰为
函数奇偶分解的精确类比，
（见\cref{ex:b1:vspaces:evenodd}），转置起着
$x \mapsto -x$ 的作用。维数方面，对称矩阵自由取值于对角线及其上方，
反对称矩阵则仅在严格上方自由取值
（对角线为零）：
\[
\frac{n(n+1)}{2} + \frac{n(n-1)}{2} = n^2 ,
\]
维数的平衡是格拉斯曼对
直和性的验证。当 $n = 2$ 时：$\begin{pmatrix} 1 & 5\\ 1 &
2\end{pmatrix} = \begin{pmatrix} 1 & 3\\ 3 & 2\end{pmatrix} +
\begin{pmatrix} 0 & 2\\ -2 & 0\end{pmatrix}$。对称
矩阵又重新出现，作为
\cref{ch:b1:multivar}中的二阶导数数据（蒙日三元组 $r, s, t$），而
对称正交矩阵则在
\cref{exo:b1:euclid:12}中被分类。
\end{example}

\section{基变换}

\begin{definition}\label{def:b1:matrices:changeofbasis}
设 $\mathcal{B}, \mathcal{B}'$ 为 $E$ 的两组基。
称 $P =
P_{\mathcal{B}\to\mathcal{B}'}$ 为\emph{过渡矩阵}\index{过渡矩阵}，其列为
\emph{新}基向量在\emph{旧}基下的坐标。
它可逆，$P^{-1} = P_{\mathcal{B}'\to\mathcal{B}}$，且
坐标按 $X = PX'$ 变换（旧 $=$ $P\,\cdot$ 新）。
\end{definition}

\begin{example}[读懂过渡矩阵]
\label{ex:b1:matrices:changerun}
在 $\R^2$ 中，从标准基 $\mathcal B$ 到 $\mathcal B' =
\bigl((1,1), (1,-1)\bigr)$：
\[
P = P_{\mathcal B\to\mathcal B'} =
\begin{pmatrix} 1 & 1\\ 1 & -1 \end{pmatrix}
\]
（新向量用旧坐标写出，逐列排列）。对于旧坐标为
$X = (3, 1)^{\mathsf T}$ 的向量，其新坐标
为 $X' = P^{-1}X = \frac12(3 + 1,\ 3 - 1)^{\mathsf T} =
(2, 1)^{\mathsf T}$。事实上 $2(1,1) + 1(1,-1) = (3,1)$。注意
方向 --- 矩阵 $P$ 由\emph{新}基构造，
但它完成的是\emph{由新到旧}的坐标换算（$X = PX'$）；由旧到新
则需取逆。写出验算
$2(1,1) + (1,-1) = (3,1)$，便能及时发现
颠倒 $P$ 这一错误，正是本章中
最常见的错误。
\end{example}

\begin{theorem}[映射的基变换]\label{thm:b1:matrices:conjugation}
设 $u \in \mathcal{L}(E)$ 的矩阵为 $A$（基为 $\mathcal{B}$），且 $A'$
为其在 $\mathcal{B}'$ 下的矩阵，且 $P = P_{\mathcal{B}\to\mathcal{B}'}$，则
\[
A' = P^{-1} A\, P .
\]
如此相关的两个矩阵称为\emph{相似矩阵}\index{相似矩阵}。
（对于 $u \colon E \to F$ 及两对基，公式为
$A' = Q^{-1} A P$ --- 即\emph{等价}矩阵
\index{等价矩阵}。）
\end{theorem}

\begin{proof}
对任意 $x$：$X = PX'$，且像满足 $Y = AX$，$Y = PY'$。
于是 $PY' = APX'$，即 $Y' = (P^{-1}AP)X'$ 对所有 $X'$ 成立，故
$u$ 在新基下的矩阵为 $P^{-1}AP$（$X'$ 取
标准列向量）。
\end{proof}

\begin{example}[好的基使映射一目了然]
\label{ex:b1:matrices:conjugationrun}
设 $u(x, y) = (y, x)$（交换），其矩阵 $A = \begin{pmatrix} 0
& 1\\ 1 & 0\end{pmatrix}$（标准基下）。在
基 $\mathcal B' = \bigl((1,1), (1,-1)\bigr)$ 下：
\[
P = \begin{pmatrix} 1 & 1\\ 1 & -1 \end{pmatrix},
\qquad
P^{-1} = \frac12\begin{pmatrix} 1 & 1\\ 1 & -1 \end{pmatrix},
\qquad
P^{-1} A P = \begin{pmatrix} 1 & 0\\ 0 & -1 \end{pmatrix}.
\]
其实并不需要真的计算矩阵乘积：$u$ 保持 $(1,1)$ 不变，并
将 $(1,-1)$ 反向，故在 $\mathcal B'$ 下其矩阵\emph{必为}
$\operatorname{diag}(1, -1)$ --- 交换恰为
关于直线 $y = x$ 的反射。对给定的自同态，寻找使其矩阵
成为对角矩阵的基，正是
本科二年级分册（约化理论）的中心问题；下文
周末问题表明，仅凭多项式恒等式本身已能走多远。
\end{example}

\begin{example}[基变换，逆向运用]
\label{ex:b1:matrices:reverserun}
到 $F = \operatorname{Vect}(1,1)$ 上、沿 $G =
\operatorname{Vect}(1,-1)$ 的投影，在适配基 $\mathcal
B' = \bigl((1,1),(1,-1)\bigr)$ 下，其矩阵为 $A' =
\operatorname{diag}(1, 0)$，形式极简。求其标准基下的矩阵，
只需将 \cref{thm:b1:matrices:conjugation} 反向运用，$A = P A'
P^{-1}$：
\[
P = \begin{pmatrix} 1 & 1\\ 1 & -1\end{pmatrix},
\quad
P^{-1} = \frac12\begin{pmatrix} 1 & 1\\ 1 & -1\end{pmatrix},
\quad
A = P\begin{pmatrix} 1 & 0\\ 0 & 0\end{pmatrix}P^{-1}
= \frac12\begin{pmatrix} 1 & 1\\ 1 & 1\end{pmatrix}.
\]
验算：$A^2 = A$（幂等），$\operatorname{tr} A = 1 =
\operatorname{rk} A$，且 $A\binom{1}{1} = \binom11$，
$A\binom{1}{-1} = 0$，与所要求的相符。这种逆向做法 ---
即先在适配基下设计矩阵，再共轭还原 ---
正是实际中构造旋转、反射与投影矩阵的
方法。
\end{example}

\begin{theorem}[秩的标准形]\label{thm:b1:matrices:rank}
矩阵的\emph{秩}（其列的秩，等价地，
相伴线性映射的秩）是等价关系下唯一的不变量：
每个 $A \in \mathcal{M}_{n,p}$，秩为 $r$，都等价于
\[
J_r = \begin{pmatrix} I_r & 0 \\ 0 & 0 \end{pmatrix},
\]
且 $\operatorname{rk}(A^{\mathsf T}) = \operatorname{rk}(A)$，即行秩
等于列秩。
\end{theorem}

\begin{proof}
设 $u \colon E \to F$ 的秩为 $r$。取 $S$ 作为
$\ker u$ 的补子空间（$\dim S = r$，见\cref{thm:b1:linmaps:ranknullity}）
其基为 $(e_1, \dots, e_r)$，再用 $\ker u$ 的基将其扩充为
$E$ 的基；像 $f_i = u(e_i)$（$i \leq r$）构成
$\operatorname{im} u$ 的基（限制映射为同构），再扩充为
$F$ 的基。在这些基下，$u$ 的矩阵恰为
$J_r$。于是 $A = Q J_r P^{-1}$，其中 $P, Q$ 可逆。

转置可得：$A^{\mathsf T} = (P^{-1})^{\mathsf T} J_r^{\mathsf T} Q^{\mathsf T}$
其中 $J_r^{\mathsf T}$ 具有相同形式（秩为 $r$），两侧的
因子可逆（可逆矩阵的转置仍可逆，由
$(AB)^{\mathsf T} = B^{\mathsf T}A^{\mathsf T}$ 应用于 $AA^{-1} = I$ 得出）：
故 $\operatorname{rk} A^{\mathsf T} = r$。
\end{proof}

\section{行变换}

\begin{method}[矩阵的高斯消元法]\label{met:b1:matrices:gauss}
三种\emph{初等行变换}\index{初等行变换} ---
交换两行、以 $\lambda \neq 0$ 乘某一行、将某行的倍数
加到另一行上 --- 都不改变秩（每种变换都是
左乘一个可逆矩阵）。算法：取主元
（最左的非零元），消去其下方同列元素，转到下一行和下一列；
所得阶梯形的主元个数
即为秩。

\emph{求逆：}对分块 $(A \mid
I_n)$ 施行算法，直至左块变为 $I_n$（当且仅当 $A$
可逆时才可能）；此时右块即为 $A^{-1}$ --- 事实上，
所用初等矩阵的乘积恰为 $A^{-1}$。
\end{method}

\begin{example}\label{ex:b1:matrices:inverse}
$A = \begin{pmatrix} 1 & 2 \\ 3 & 4 \end{pmatrix}$：对 $(A \mid
I_2)$ 作行化简：
\[
\begin{pmatrix} 1 & 2 & 1 & 0\\ 3 & 4 & 0 & 1 \end{pmatrix}
\to
\begin{pmatrix} 1 & 2 & 1 & 0\\ 0 & -2 & -3 & 1 \end{pmatrix}
\to
\begin{pmatrix} 1 & 0 & -2 & 1\\ 0 & 1 & \tfrac32 & -\tfrac12
\end{pmatrix},
\]
（运算：$L_2 \leftarrow L_2 - 3L_1$；再令 $L_1 \leftarrow L_1
+ L_2$，$L_2 \leftarrow -\frac12 L_2$）。于是 $A^{-1} =
\begin{pmatrix} -2 & 1 \\ \tfrac32 & -\tfrac12\end{pmatrix}$.
\emph{验算：}$AA^{-1} = I_2$。
\end{example}

\begin{example}[含参数的秩，仅用行变换]
\label{ex:b1:matrices:rankparameter}
设 $m \in \R$，求
$M_m = \begin{pmatrix} 1 & 1 & m\\ 1 & m & 1\\ m & 1 &
1\end{pmatrix}$ 的秩。作行化简：$L_2 \leftarrow L_2 - L_1$ 与 $L_3
\leftarrow L_3 - mL_1$ 后，各行分别为
\[
(1,\ 1,\ m), \qquad (0,\ m - 1,\ 1 - m), \qquad
(0,\ 1 - m,\ 1 - m^2).
\]
\emph{情形 $m = 1$}：后两行化为零，只有一个主元，
故 $\operatorname{rk} M_1 = 1$（原三行
全相等）。\emph{情形 $m \neq 1$}：将 $L_2$ 乘以 $\frac1{m-1}$，
将 $L_3$ 乘以 $\frac1{1-m}$，得 $(0, 1, -1)$ 与 $(0, 1, 1 + m)$，
再令 $L_3 \leftarrow L_3 - L_2 = (0, 0, m + 2)$。若 $m = -2$，
则有两个主元，秩为 $2$；否则有三个主元，秩为 $3$。综上：
\[
\operatorname{rk} M_m =
\begin{cases}
1 & m = 1,\\
2 & m = -2,\\
3 & \text{否则}.
\end{cases}
\]
同样的分界点也会从一次行列式计算中得出，
在\cref{ch:b1:det}中，多项式 $-(m+2)(m-1)^2$（
见\cref{exo:b1:det:7}）---但要注意，消元法所给出、
而行列式所不给出的，是退化情形下秩的\emph{值}，
而不仅是秩下降这一事实。
\end{example}

\begin{example}[计算幂]\label{ex:b1:matrices:powers}
$A = \begin{pmatrix} 1 & 1 \\ 0 & 1\end{pmatrix} = I + N$，其中 $N =
E_{12}$，$N^2 = 0$。由于 $I$ 与 $N$ 可交换，二项式定理
（\cref{prop:b1:structures:binomial}）化为有限和：
\[
A^k = I + kN = \begin{pmatrix} 1 & k \\ 0 & 1 \end{pmatrix}
\qquad (k \in \N, \text{ 且 } k \in \Z \text{ 利用 } A^{-1} = I -
N).
\]
\end{example}

\begin{method}[计算 $A^n$：三种途径]
\label{met:b1:matrices:powers}
\begin{enumerate}
  \item \emph{二项式途径}：若 $A = \lambda I + N$，其中 $N$
        幂零，二项式定理化为有限和
        （\cref{ex:b1:matrices:powers}，\cref{exo:b1:matrices:5}）；
        此途径适用，因为 $\lambda I$ 与一切矩阵可交换。
  \item \emph{多项式途径}：找一个多项式恒等式，
        被 $A$ 满足（维数为 $2$ 时恒有 $A^2 = sA -
        pI$），并将 $X^n$ 模之约化；
        下文的周末问题完整地建立了这一途径。
  \item \emph{相似途径}：找可逆矩阵 $P$，使
        $P^{-1}AP = D$ 为简单形（对角阵或移位阵），计算
        $D^n$，再还原：$A^n = P D^n P^{-1}$
        （\cref{thm:b1:matrices:conjugation}，
        \cref{ex:b1:matrices:conjugationrun}）；寻找这类 $P$ 的系统方法
        属于本科二年级的约化理论。
\end{enumerate}
无论采用哪一种途径，都应在 $n = 0, 1, 2$ 时检验结果：三项
简便的检验，几乎能查出所有失误。
\end{method}

\begin{remark}[常见陷阱：不可交换性的代价]
\label{rem:b1:matrices:pitfalls}
凡数量代数中成立、其证明须交换因子次序的恒等式，
在 $\mathcal{M}_n(K)$，$n \geq 2$ 时均失效。
\emph{平方}：$(A + B)^2 = A^2 + AB + BA + B^2$，中间项
并为 $2AB$，仅当 $AB = BA$
（见\cref{exo:b1:matrices:1}）。
\emph{积的幂}：$(AB)^k$ 为 $ABAB\cdots$，而非
$A^kB^k$。
\emph{零因子}：$E_{12}E_{12} = 0$，而 $E_{12} \neq 0$；
因此 \emph{不满足消去律}：$AB = AC$ 推出 $B = C$
仅当 $A$ 可逆时才成立（在正确一侧乘 $A^{-1}$ ---
注意左右之分）。
\emph{迹}：$\operatorname{tr}(AB) = \operatorname{tr}(BA)$
恒成立，但 $\operatorname{tr}(AB) \neq
\operatorname{tr}A\operatorname{tr}B$ 一般不成立（取 $A = B =
I_2$：$2 \neq 4$），且 $\operatorname{tr}(ABC) =
\operatorname{tr}(BCA)$（轮换），而 $\operatorname{tr}(ACB)$
可能不同。
\emph{转置反序}：$(AB)^{\mathsf T} = B^{\mathsf
T}A^{\mathsf T}$ --- 忘记反序是正交性计算中最常见的
错误（见\cref{ch:b1:euclid}）。
拿不准时，可取 $E_{12}$ 与
$E_{21}$ 代入检验任一所称恒等式：最小的不可交换一对，
一行之内便可推翻大多数错误公式。
\end{remark}

\begin{remark}[矩阵字典的去向]
\label{rem:b1:matrices:whereused}
矩阵字典在本册余下的每一部分都会用到：
\cref{ch:b1:det}为每个方阵赋予一个
决定可逆性的数，并系统地求解 $AX = B$；
\cref{ch:b1:euclid}单独挑出保持长度的矩阵
（即正交矩阵）；并在
\cref{ch:b1:multivar}中，二元函数的二阶性态是
一个对称的 $2 \times 2$ 矩阵。前文几乎一笔带过的迹，
则成为一个强有力的不变量：
\cref{exo:b1:matrices:6,exo:b1:matrices:8}给出初步体会，
本科二年级分册则在此基础上建立特征值理论。
周末问题发展了另一件主力工具：多项式恒等式
矩阵所满足的，它们把 $A^n$ 的计算化为
一个两项线性递推式。
\end{remark}

\begin{remark}[第 3 册内的展望]
\label{rem:b1:matrices:perspectives}
这里引入的三类矩阵，将在
本册后文中再次出场。\emph{对称矩阵}
（\cref{ex:b1:matrices:symsplit}）承载
二元函数的二阶信息：
\cref{ch:b1:multivar}中的蒙日判别法是关于
一个对称 $2\times2$ 矩阵符号性态的论断，其行列式 $rt - s^2$
由\cref{ch:b1:det}中的方法算出。\emph{正交
矩阵}（$A^{\mathsf T}A = I$）是
\cref{ch:b1:euclid}中的等距变换，转置在那里终于获得了
几何意义：它是内积的
代数影子。\emph{可逆矩阵}在
\cref{ch:b1:det}中迎来实用检验 --- 只需一个数 $\det A \neq 0$ ---
为本章开始的行化简探寻画上句号。
迹与行列式随后作为周末问题中的不变量对 $(s, p)$ 结伴出场，
一直延伸到
本科二年级的特征值理论。
\end{remark}

\section{练习}

\begin{exercise}[$\star$]\label{exo:b1:matrices:1}
设 $A = \begin{pmatrix} 1 & 2 \\ 0 & 1 \end{pmatrix}$ 与 $B =
\begin{pmatrix} 0 & 1 \\ 1 & 0\end{pmatrix}$。计算 $AB$、$BA$、
$A^2 - B^2$ 与 $(A+B)(A-B)$；并说明后两者为何不相等。
\end{exercise}

\begin{exercise}[$\star$]\label{exo:b1:matrices:2}
求下列矩阵的秩：
\[
M = \begin{pmatrix}
1 & 2 & 3\\
2 & 4 & 6\\
1 & 1 & 1
\end{pmatrix},
\qquad
N = \begin{pmatrix}
1 & 1 & 0 & 2\\
0 & 1 & 1 & 1\\
1 & 2 & 1 & 3
\end{pmatrix}.
\]
\end{exercise}

\begin{exercise}[$\star$]\label{exo:b1:matrices:3}
用初等行变换求 $A = \begin{pmatrix} 1 & 0 & 1\\ 2 & 1 &
1\\ 1 & 1 & 1 \end{pmatrix}$ 的逆矩阵，并用一次乘积验证。
\end{exercise}

\begin{exercise}[$\star$]\label{exo:b1:matrices:4}
在 $\R_2[X]$ 的标准基下，写出自同态
$u(P) = P(X + 1)$ 的矩阵，不经计算说明其为何可逆，
并给出 $u^{-1}$ 的矩阵。
\end{exercise}

\begin{exercise}[$\star\star$]\label{exo:b1:matrices:5}
设 $A = \begin{pmatrix} 2 & 1 \\ 0 & 2\end{pmatrix}$，写出 $A =
2I + N$，计算 $N^2$，并
由二项式定理推出 $A^k$，对所有 $k \in \N$ 成立。
\end{exercise}

\begin{exercise}[$\star\star$]\label{exo:b1:matrices:6}
证明：不存在矩阵 $A, B \in \mathcal{M}_n(K)$（$K = \R$ 或 $\C$），
使得 $AB - BA = I_n$。\emph{（取迹。）}
\end{exercise}

\begin{exercise}[$\star\star$]\label{exo:b1:matrices:7}
称矩阵 $A$ 为\emph{幂零矩阵}，若 $A^m = 0$ 对某个 $m$ 成立。
证明：此时 $I - A$ 可逆，且
\[
(I - A)^{-1} = I + A + A^2 + \dots + A^{m-1} .
\]
应用：求矩阵 $\begin{pmatrix} 1 & 2 & 3\\ 0 & 1 & 2\\ 0 & 0 &
1\end{pmatrix}$ 的逆矩阵。
\end{exercise}

\begin{exercise}[$\star\star$]\label{exo:b1:matrices:8}
设 $A \in \mathcal{M}_n(\R)$ 满足 $A^2 = A$（即幂等矩阵）。证明
$\operatorname{tr} A = \operatorname{rk} A$。\emph{（把
$A$ 视为投影，并选取适配的基；
\cref{thm:b1:matrices:conjugation} 表明迹
与基的选取无关，因为 $\operatorname{tr}(P^{-1}MP) =
\operatorname{tr} M$。）}
\end{exercise}

\begin{exercise}[$\star\star\star$]\label{exo:b1:matrices:9}
设 $J \in \mathcal{M}_n(\R)$ 为全一矩阵。计算 $J^2$，
并由此推出，对 $a, b \in \R$，矩阵 $M
= aI + bJ$ 可逆的条件及 $M^{-1}$ \emph{（试求同形式
$\alpha I + \beta J$ 的逆）}。
\end{exercise}

\begin{exercise}[$\star\star\star$]\label{exo:b1:matrices:10}
（秩不等式）对 $A, B \in \mathcal{M}_n(K)$，证明
\[
\operatorname{rk}(A + B) \leq \operatorname{rk} A +
\operatorname{rk} B,
\qquad
\operatorname{rk}(AB) \geq \operatorname{rk} A + \operatorname{rk}
B - n .
\]
\emph{（对第二式 --- 西尔维斯特不等式 --- 将秩–零化度定理
应用于 $A$ 所对应的线性映射在
$\operatorname{im} B$ 上的限制。）}
\end{exercise}

\begin{exercise}[$\star\star$]\label{exo:b1:matrices:11}
设 $D = \operatorname{diag}(d_1, \dots, d_n)$，其中 $d_i$
\emph{两两互异}。
\begin{enumerate}
  \item 证明矩阵 $A$ 与 $D$ 可交换，当且仅当
        $A$ 为对角矩阵。\emph{（比较 $(i,j)$ 项，
        即 $AD$ 与 $DA$ 的对应元素。）}
  \item 由此推出 $\mathcal{M}_n(K)$ 的\emph{中心}：
        与\emph{任一}矩阵可交换的矩阵恰为
        数量矩阵 $\lambda I_n$。\emph{（先
        以 $D$ 作检验，再以矩阵 $E_{ij}$ 作检验。）}
\end{enumerate}
\end{exercise}

\begin{exercise}[$\star\star\star$]\label{exo:b1:matrices:12}
（秩为 1 的矩阵）设 $A \in \mathcal{M}_n(K)$，$A \neq 0$。
\begin{enumerate}
  \item 证明：$\operatorname{rk} A = 1$ 当且仅当 $A =
        CL$，其中 $C \in \mathcal{M}_{n,1}$ 为非零列，
        $L \in \mathcal{M}_{1,n}$ 为非零行。
  \item 对这样的 $A$，证明 $A^2 = (\operatorname{tr} A)\,A$；
        由此推出：秩为 1 的矩阵幂零，当且仅当
        其迹为零。
  \item 若 $\operatorname{tr} A \neq -1$，证明 $I_n + A$
        可逆，且
        \[
        (I_n + A)^{-1} = I_n - \frac{1}{1 + \operatorname{tr}
        A}\,A ,
        \]
        且 $I_n + A$ 在
        $\operatorname{tr} A = -1$ 时\emph{不}可逆。\emph{（求一个被
        $I_n + A$ 零化的向量。）}
\end{enumerate}
\end{exercise}

\section{问题：借助多项式除法求矩阵的幂}

\begin{problem}\label{pb:b1:matrices:1}
逐个元素计算 $A^{100}$ 毫无希望，
而借助 $A$ 所满足的多项式恒等式计算，只需三行。
本题从头建立这一方法：$X^n$ 的带余除法，
恒等式 $A^2 - sA + pI = 0$ 对每个 $2
\times 2$ 矩阵都成立（即
$2$ 维时的凯莱–哈密顿定理），并给出矩阵幂与线性
递推之间的对应关系 --- 以斐波那契数为贯穿全题的例子。
\index{凯莱–哈密顿定理（2 维）}
\index{斐波那契数}\index{线性递推}

\medskip
\noindent\textbf{第一部分 --- 余式演算。} 设 $s, p
\in K$ 且 $D = X^2 - sX + p$。
\begin{enumerate}
  \item 说明对每个 $n \in \N$，存在唯一的 $Q_n
        \in K[X]$ 与 $(a_n, b_n) \in K^2$，使得
        \[
        X^n = Q_n\,D + a_n X + b_n ,
        \]
        并计算 $(a_0, b_0)$ 与 $(a_1, b_1)$。
  \item 两边乘以 $X$ 再作一次除法，建立
        下列递推关系
        \[
        a_{n+1} = s\,a_n + b_n,
        \qquad
        b_{n+1} = -p\,a_n ,
        \]
        并由此推出 $a_{n+2} = s\,a_{n+1} - p\,a_n$：
        系数数列满足与
        $D$ 相应的线性递推式。
  \item 设 $D$ 有两个不同的根 $\lambda \neq \mu$。
        在带余除法恒等式中代入根，证明
        \[
        a_n = \frac{\lambda^n - \mu^n}{\lambda - \mu},
        \qquad
        b_n = \frac{\lambda\mu^n - \mu\lambda^n}{\lambda - \mu} .
        \]
  \item 设 $D = (X - \lambda)^2$。对除法恒等式求导，
        证明 $a_n = n\lambda^{n-1}$ 与
        $b_n = (1 - n)\lambda^{n}$。
  \item 证明将固定矩阵 $M \in
        \mathcal{M}_k(K)$ 代入多项式时保持和与积：
        $(PQ)(M) = P(M)\,Q(M)$。由此推出：若 $D(M) =
        0$，则
        \[
        M^n = a_n\,M + b_n\,I \qquad (n \in \N).
        \]
\end{enumerate}

\medskip
\noindent\textbf{第二部分 --- 二维情形：迹、
行列式值、凯莱–哈密顿。} 设 $A = \begin{pmatrix} a & b\\ c &
d\end{pmatrix}$ 记 $s = a + d = \operatorname{tr} A$，$p = ad
- bc$（即在\cref{ch:b1:det}中将称为
行列式的数）。
\begin{enumerate}[resume]
  \item 直接计算验证\emph{$2$ 维凯莱–哈密顿
        恒等式}：
        \[
        A^2 - s\,A + p\,I_2 = 0 .
        \]
  \item 用直接展开证明 $p$ 具有乘性：
        记号意义明确，即 $p(AB) = p(A)\,p(B)$。再证明：
        $A$ 可逆当且仅当 $p \neq 0$，此时
        有
        \[
        A^{-1} = \frac1p\,\bigl(s\,I_2 - A\bigr).
        \]
  \item 设 $A = \begin{pmatrix} 1 & 1\\ 0 & 2\end{pmatrix}$。
        计算 $s$、$p$ 与 $D$ 的根，并推出 $A^n$ 的
        闭式，再与直接计算的
        $A^2$ 进行检验。
  \item 设 $A = \begin{pmatrix} 3 & 1\\ -1 & 1\end{pmatrix}$。
        证明 $D$ 有二重根，并计算 $A^n$；
        取 $n = 2$ 进行检验。
  \item 设 $F = \begin{pmatrix} 1 & 1\\ 1 & 0\end{pmatrix}$，
        由 $F_0 = 0$，$F_1 = 1$，
        $F_{n+2} = F_{n+1} + F_n$ 定义斐波那契数。证明
        \[
        F^n = \begin{pmatrix} F_{n+1} & F_n\\ F_n &
        F_{n-1}\end{pmatrix} \quad (n \geq 1),
        \]
        由此推出比内公式 $F_n = \dfrac{\varphi^n -
        \psi^n}{\sqrt5}$，其中 $\varphi = \frac{1 + \sqrt5}2$，
        $\psi = \frac{1 - \sqrt5}2$，并利用第 7 题，
        推出卡西尼恒等式 $F_{n+1}F_{n-1} - F_n^2 = (-1)^n$。
\end{enumerate}

\medskip
\noindent\textbf{第三部分 --- 线性递推的结构。}
设 $s, p \in K$ 且 $p \neq 0$，记 $E_D$ 为
满足 $u_{n+2} = s\,u_{n+1} - p\,u_n$ 对一切 $n$ 成立的数列全体。
\begin{enumerate}[resume]
  \item 证明 $E_D$ 是维数为 $2$ 的向量空间
        （仿照\cref{exo:b1:findim:10}）。
  \item 证明第一部分中的数列 $(a_n)$ 是
        $E_D$ 中满足初值 $0, 1$ 的元素，并证明每个 $u
        \in E_D$ 都满足
        \[
        u_n = u_1\,a_n + u_0\,b_n \qquad (n \in \N),
        \]
        其中 $(b_n)$ 同第一部分：带余除法的余式
        一次性解出\emph{全部}递推式。
  \item 若 $\lambda \neq \mu$ 为 $D$ 的两个根，证明
        $\bigl((\lambda^n), (\mu^n)\bigr)$ 是 $E_D$ 的一组基；
        若 $D = (X-\lambda)^2$ 且 $\lambda \neq 0$，证明
        $\bigl((\lambda^n), (n\lambda^n)\bigr)$ 亦为一组基。
  \item 完全求解：$u_{n+2} = u_{n+1} + 6u_n$，$u_0 = 1$，
        $u_1 = 8$；用 $u_2$ 与 $u_3$ 检验结果。
  \item 设 $C = \begin{pmatrix} 0 & 1\\ -p & s\end{pmatrix}$
        （$D$ 的\emph{友矩阵}）。证明
        \[
        \begin{pmatrix} u_{n}\\ u_{n+1}\end{pmatrix}
        = C^n \begin{pmatrix} u_0\\ u_1\end{pmatrix}
        \quad (u \in E_D),
        \]
        并证明 $\operatorname{tr} C = s$ 与 $p(C) = p$：
        递推式与矩阵共同对应同一个多项式 $D$。
\end{enumerate}

\medskip
\noindent\textbf{第 IV 部分 --- 三次情形。}设 $D_3 = X^3 -
\alpha X^2 - \beta X - \gamma$，并令
\[
C_3 = \begin{pmatrix} 0 & 1 & 0\\ 0 & 0 & 1\\ \gamma & \beta &
\alpha \end{pmatrix}.
\]
\begin{enumerate}[resume]
  \item 证明 $D_3(C_3) = 0$。\emph{（计算标准基向量在
        $C_3$ 各次幂作用下的像：
        $C_3$ 的映射将 $e_1 \mapsto \dots \mapsto$ 化为
        由最后一行所确定的线性组合。）}
  \item 证明：若 $D_3$ 有三个互异的根 $\lambda_1,
        \lambda_2, \lambda_3$，则 $R_n$ 为 $X^n$
        除以 $D_3$ 所得的余式，是\emph{拉格朗日插值多项式}，
        它以 $\lambda_i^n$ 为取值、$\lambda_i$ 为节点
        （见\cref{thm:b1:poly:lagrange}）；由此可知
        $C_3^{\,n}$ 的每个元素都是
        $\lambda_1^n, \lambda_2^n, \lambda_3^n$ 的固定线性组合。
  \item 解递推式 $u_{n+3} = 2u_{n+2} + u_{n+1} - 2u_n$，其中 $u_0 =
        0$，$u_1 = 1$，$u_2 = 1$。\emph{（分解 $D_3 = (X - 1)(X
        + 1)(X - 2)$。）}用 $u_3$ 检验。
  \item 求 $X^n$ 模 $(X -
        \lambda)^3$ 的余式\emph{（$X^n$ 在
        $\lambda$ 处的泰勒展开）}，并由此推出 $(\lambda I +
        N)^n$ 的表达式，其中 $N^3 = 0$ 且 $N$ 与所有相关对象可交换；
        用二项式定理检验此式。
  \item 证明：对于有互异根的 $D_3$，$3$ 阶递推式的通解
        为 $u_n = c_1
        \lambda_1^n + c_2\lambda_2^n + c_3\lambda_3^n$，并证明
        这三个等比数列构成解空间的
        一组基。\emph{（证明线性无关时，对零组合
        在 $n = 0, 1, 2$ 处求值，并识别出
        它是关于互异节点的插值方程组，
        节点为 $\lambda_i$。）}
\end{enumerate}

\medskip
\noindent\textbf{第 V 部分 --- 斐波那契数的推论与综合。}
\begin{enumerate}[resume]
  \item 证明 $F_1 + F_2 + \dots + F_n = F_{n+2} - 1$。
  \item 由 $F^{m+n} = F^m F^n$ 推出加法公式
        \[
        F_{m+n} = F_{m+1}F_n + F_m F_{n-1},
        \]
        并由此推出 $F_{2n} = F_n(F_{n+1} + F_{n-1})$。
  \item 证明 $F_n$ 为最接近
        $\varphi^n/\sqrt5$ 的整数（对每个 $n \geq 0$）。
  \item 记 $t_n = \operatorname{tr}(F^n) = F_{n+1} + F_{n-1}$
        （即\emph{卢卡斯数} $L_n$）。证明 $t_{n+2} =
        t_{n+1} + t_n$，$t_1 = 1$，$t_2 = 3$；并证 $L_n =
        \varphi^n + \psi^n$，进而推出 $F_{2n} = F_n L_n$。
  \item 综合，用四句话回答：为何 $2
        \times 2$ 矩阵的幂位于平面
        $\operatorname{Vect}(I, A)$（$\mathcal{M}_2(K)$ 的子空间），
        即哪一个维数论证保证了二次恒等式，
        以及第 II 部分给出了哪个显式恒等式；欧几里得除法如何
        把幂运算化为两项递推式；
        本题中哪一个命题是在所有维数下都成立的某定理当 $n =
        2$ 时的特例（写出其名称，
        并说明它在本系列中于何处被证明）；以及
        友矩阵构造为图景增添了什么。
\end{enumerate}
\end{problem}
